\documentclass[t]{beamer}

\usetheme{CambridgeUS}

\title{IM1003: Programming Design, Spring 2017  \linebreak Lab 01}
\author[bigelephant29]{Jhih-Bang Hsieh\linebreak \small{bigelephant29}}
\institute{\textbf{National Taiwan University}}
\date{}

\usefonttheme{serif}
\usepackage{xeCJK} 
\usepackage{fontspec}
\setCJKmainfont{Adobe Ming Std}

\usepackage{graphicx}
\graphicspath{{/}}

\makeatletter
\setbeamertemplate{footline}{%
    \leavevmode%
    \hbox{%
        \begin{beamercolorbox}[wd=.8\paperwidth,ht=2.25ex,dp=1ex,center]{title in head/foot}%
            \usebeamerfont{title in head/foot}\insertshorttitle
        \end{beamercolorbox}%
        \begin{beamercolorbox}[wd=.2\paperwidth,ht=2.25ex,dp=1ex,right]{date in head/foot}%
            \usebeamerfont{date in head/foot}\insertshortdate{}\hspace*{2em}
            \insertframenumber{} / \inserttotalframenumber\hspace*{2ex} 
        \end{beamercolorbox}}%
        \vskip0pt%
    }
\makeatother

\begin{document}

% First Page
\begin{frame}
  \maketitle
\end{frame}

% Outline Page
\begin{frame}{Outline}
  \begin{itemize}
    \item TA \& Lab
    \item Introduction to PDOGS
    \item IDE
    \item Text Editor
    \item Installing g++
    \item Coding Style
    \item Practice
  \end{itemize}
\end{frame}

% Section: TA & Lab
\section{TA \& Lab}
\begin{frame}{TA}
  \begin{itemize}
    \setlength\itemsep{1em}
    \item
    李維哲\ Jeff Lee\\
    r04725023(AT)ntu.edu.tw
    \item
    林敬傑\ Jack Lin\\
    r05725007(AT)ntu.edu.tw
    \item
    楊佩蓉\ Sophie Yang\\
    r05725028(AT)ntu.edu.tw
    \item
    謝志邦\ Jhih-Bang Hsieh\\
    b02705021(AT)ntu.edu.tw
  \end{itemize}
\end{frame}

\begin{frame}{Lab}
  \begin{itemize}
    \item 地點：管理學院一號館\ 三樓大電腦教室
    \item 時間：星期三 18:30 - 20:15
    \item Lab 不計入學期成績，可自由參加。
    \item Lab 的練習題應該會很有趣。
    \item 四位助教每週 Lab 都會在，不另開 Office hour。
    \item Lab 的所有資料以及練習題參考程式碼皆會放在助教的 Github。
    \item 有任何錯誤歡迎指正！
    \item Link: \href{https://github.com/bigelephant29/PD17-Lab}{https://github.com/bigelephant29/PD17-Lab}
  \end{itemize}
\end{frame}

% Section: Introduction to PDOGS
\section{Introduction to PDOGS}
\begin{frame}{Introduction to PDOGS}
  \begin{itemize}
    \item Programming Design Online Grading System
    \item 由資管系 B02、B03、B04 同學齊心協力開發。
    \item 可以用 @ntu.edu.tw 信箱申請帳號。
    \item 系統應該是 24 小時活著，發現系統不小心死掉請戳戳四位獸醫。
    \item 作業、考試皆在系統上，系統備有自動評分功能。
    \item 請不要用任何方式攻擊系統（包含物理攻擊），不然這門課你可能要修兩次以上。
  \end{itemize}
  \begin{center}
    \includegraphics[width=5em]{pdogs.png}
  \end{center}
\end{frame}

% Section: IDE
\section{IDE}
\begin{frame}{IDE}
  \begin{itemize}
    \item Integrated Development Environment 整合開發環境
    \item 整合了文字編輯器（Text Editor）、編譯器（Compiler）的環境。
    \item 有時還會包含除錯器（Debugger）等輔助開發工具。
    \item 常見的 IDE：Dev-C++, Code::Blocks, Xcode, Eclipse ...等。
  \end{itemize}
  \begin{center}
    \includegraphics[height=2.5em]{dev.png}
    \hspace{1em}
    \includegraphics[height=2.5em]{cb.png}
    \hspace{1em}
    \includegraphics[height=2.5em]{xcode.png}
    \hspace{1em}
    \includegraphics[height=2.5em]{eclipse.png}
  \end{center}
\end{frame}

\begin{frame}{Dev-C++}
  \includegraphics[height=2.5em]{dev.png}
  \begin{itemize}
    \item 對初學者非常友善的一個 IDE。
  \end{itemize}
\end{frame}

% Section: Text Editor
\section{Text Editor}
\begin{frame}{Text Editor}
\end{frame}

% Section: Installing g++
\section{Installing g++}
\begin{frame}{Installing g++}
\end{frame}

% Section: Coding Style
\section{Coding Style}
\begin{frame}{Coding Style}
\end{frame}

% Section: Practice
\section{Practice}
\begin{frame}{Practice}
\end{frame}

\end{document}
